\documentclass[11pt]{article}
\usepackage[margin=0.9in]{geometry}
\usepackage[T1]{fontenc}
\usepackage{lmodern,amsmath,amssymb,amsthm,microtype}
\usepackage[colorlinks=true,urlcolor=blue,citecolor=blue,linkcolor=blue]{hyperref}
\theoremstyle{definition}
\newtheorem{theorem}{Theorem}
\newtheorem{lemma}[theorem]{Lemma}
\newcommand{\PP}{\mathbb P}
\newcommand{\ZZ}{\mathbb Z}
\newcommand{\NN}{\mathbb N}
\title{A supercritical nonuniqueness example on a graph\\
roughly isometric to the square lattice}
\author{Alec Kriebel\\ \small Independent researcher\\
\small \href{https://orcid.org/0009-0001-9320-500X}{ORCID 0009-0001-9320-500X}}
\date{22 October 2026 \quad Version 1.0.0}
\begin{document}
\maketitle
\begin{abstract}
We give an explicit deterministic graph of maximum degree four, roughly
isometric to the square lattice, on which iid bond percolation has at least
two infinite clusters with probability greater than one third at parameter
$19/20$, while its critical probability is at most $9/10$.
Edges internal to two opposite quadrants are retained; every other lattice
edge is replaced by a path of length $100$. Elementary dual parity and
all-distance path bounds prove the assertions without a lattice threshold
formula or a transitivity assumption. The example refutes the explicitly
declared all-supercritical interpretation of Benjamini's 2013 Problem 9.6.
The original question omits its probability quantifier; no unique author
intent, exclusive priority, or current open status is asserted.
\end{abstract}

\section{Statement and the source question}
All graphs here are countable, simple, undirected and connected, with graph
distance. In Bernoulli bond percolation every edge is independently retained
with the same probability $p$. Write $N_\infty$ for the number of infinite
open components and
\[
 p_c(G,o)=\inf\{p:\PP_p(o\leftrightarrow\infty)>0\}.
\]
For a connected graph, positive survival at one vertex implies positive survival
at every vertex at the same $p>0$: force open a fixed finite connecting path.
Indeed, conditioning those finitely many bonds to be open can only increase
any increasing event, by coupling the other bonds identically. Hence
$\PP_p(o\leftrightarrow\infty)\geq p^{|F|}\PP_p(v\leftrightarrow\infty)$
for such a path $F$ from $o$ to $v$. We may therefore write $p_c(G)$.

We use precisely the multiplicative-plus-additive rough-isometry convention
of~\cite[Definition 1.25, manuscript p.~9]{benjamini}: a map
$f:V(G)\to V(H)$ and a constant $C>0$ satisfy
\[
 C^{-1}d_H(f(x),f(y))-C\leq d_G(x,y)
 \leq C d_H(f(x),f(y))+C,
\]
and every vertex of $H$ is at distance strictly less than $C$ from the range.

Problem 9.6 of the 30 October 2013 author manuscript~\cite[p.~64]{benjamini}
asks for uniqueness on bounded-degree graphs roughly isometric to $\ZZ^d$.
Its sentence supplies no $p$-quantifier. The chapter's preceding page explicitly
discusses $p>p_c$, while transitivity is a separate hypothesis of nearby
theorems. We address the following \emph{declared interpretation}:
for every such graph and every fixed $p_c(G)<p\leq1$,
$\PP_p(N_\infty=1)=1$. A negative example in dimension two refutes this
interpretation. It does not refute a weaker assertion of uniqueness only
sufficiently near one, determine the author's unexpressed intent, or require
a simultaneous assertion for all $p$ in one coupling.

\begin{theorem}\label{main}
There is a deterministic graph $G$ of maximum degree four, roughly isometric
to $\ZZ^2$ with $C=100$, such that
\[
 p_c(G)\leq\frac9{10}<\frac{19}{20},
 \qquad
 \PP_{19/20}(N_\infty\geq2)>
 \frac{2200741829}{6461571200}>\frac13 .
\]
At $p=9/10$ every vertex has positive probability of an infinite open cluster.
\end{theorem}

\section{The graph and the exact coarse model}
In the nearest-neighbor square lattice let
\[
 A=\{(x,y):x\geq1,\ y\geq1\},\qquad
 B=\{(x,y):x\leq-1,\ y\leq-1\}.
\]
Call edges internal to $A$ or internal to $B$ strong; all others are weak.
Keep each strong edge, and replace every weak edge by a path of exactly
$100$ edges with $99$ new internal vertices private to that path.
There are no other edges. Original vertices have degree four; new vertices
have degree two. The resulting graph is infinite, countable, connected,
simple and deterministic. The subdivision bound is global, not a length
growing with location.

Fix original vertices under $f:G\to\ZZ^2$, and send an internal path vertex
to its nearer original endpoint, choosing ties consistently. Adjacent
vertices have equal or neighboring images, so
$d_{\ZZ^2}(f(x),f(y))\leq d_G(x,y)$. Each vertex is within $50$ actual
edges of its image, and a lattice geodesic between original vertices expands
to a path at most $100$ times as long. Thus, for all vertices,
\[
 d_{\ZZ^2}(f(x),f(y))\leq d_G(x,y)
 \leq100d_{\ZZ^2}(f(x),f(y))+100.
\]
Surjectivity of $f$ supplies the strict $100$-net condition. These inequalities
imply the displayed source convention with the single constant $100$,
including for chain midpoints.

At actual bond parameter $p$, declare a coarse strong edge open when its
single bond is open, and a coarse weak edge open when \emph{all} its
$100$ bonds are open. Coarse states are independent with probabilities
$p$ and $q=p^{100}$, respectively: their actual bond sets are disjoint.
Sharing endpoints creates no dependence in a bond model. This familiar
series reduction is used, for example, in~\cite[Section 2]{hsst};
it is not a new probabilistic mechanism.

Actual and coarse connectivity of original vertices agree. A simple actual
path between originals must traverse every chain it enters completely,
since its internal vertices have degree two. Conversely, expand open
coarse edges. An actual component containing no original lies inside one
finite chain. A component containing finitely many originals lies in
those vertices and their finitely many incident finite chains, and is
finite. Hence every infinite component contains infinitely many originals.
Restriction to originals and expansion therefore give a bijection between
actual and coarse infinite components, even when partial chain branches
are present. The inhomogeneous coarse description is exactly the requested
homogeneous iid model on $G$.

\section{Survival inside a strong quadrant}
Let $Q$ be the nearest-neighbor graph on $\NN_0^2$.
We prove the explicit bound
\begin{equation}\label{survive}
 \PP_{9/10}^Q\bigl((0,0)\leftrightarrow\infty\bigr)\geq\frac{143}{245}.
\end{equation}
For a finite set $D\subset\NN_0^2$ containing the origin, draw the dual
edges crossing primal $Q$-edges with exactly one endpoint in $D$.
Interior dual vertices are square centers $(i+1/2,j+1/2)$.
Boundary dual endpoints are $(i+1/2,-1/2)$ and $(-1/2,j+1/2)$,
where $i,j\geq0$.

At an interior dual vertex the degree is $0$, $2$ or $4$, since the cyclic
binary membership sequence around its square has an even number of changes.
A bottom or left boundary vertex has degree at most one.
On the primal bottom ray, membership starts at one and is eventually zero,
so the number of bottom degree-one endpoints is odd. Some finite dual
component consequently has an odd number of bottom endpoints. By the
handshaking lemma it also has an odd, and hence positive, number of left
endpoints. A simple path in that component joins a bottom endpoint
$(i+1/2,-1/2)$ to a left endpoint $(-1/2,j+1/2)$.
Boundary vertices cannot be internal because they have degree one.
Its length $m$ satisfies
\[
 m\geq i+j+2\geq2.
\]
The argument allows holes, disconnected $D$, diagonal contacts and degree-four
dual vertices; no choice of an outer-boundary tour is required.

If the open corner cluster is finite, take it as $D$. Every crossed primal
edge is closed. For each length $m\geq2$ there are at most $m$ possible
bottom starts, since $i\leq m-2$. From one start there are at most
$4\cdot3^{m-1}$ self-avoiding dual paths of length $m$.
Each eligible path crosses $m$ distinct actual $Q$-edges, so its probability
of being closed is $(1-t)^m$. Ineligible paths enlarge only the count;
no probability is assigned to edges outside $Q$.
For $z=3(1-t)<1$, the union bound gives
\begin{equation}\label{peierls}
 \PP_t^Q(\text{corner cluster finite})
 \leq\frac43\sum_{m\geq2}m z^m
 =\frac43\,\frac{z^2(2-z)}{(1-z)^2}.
\end{equation}
The series identity follows by subtracting its $m=1$ term from
$\sum_{m\geq1}m z^m=z/(1-z)^2$.
At $t=9/10$, the last expression is $102/245$, proving~\eqref{survive}.
In particular $p_c(G)\leq9/10$, and the finite-path forcing argument
gives positive survival at every actual vertex. Coupling edge states
by common independent uniforms makes quadrant survival increasing in $p$.
Thus at $p=19/20$ the events $S_A,S_B$ that the respective corners
$(1,1),(-1,-1)$ survive using only their internal strong edges each have
probability at least $143/245$. This uses monotonicity of survival,
not monotonicity of uniqueness.

\section{A bound on every weak connection}
Fix $p=19/20$ and put $q=p^{100}$, $r=3q$.
If an open coarse path connects $A$ to $B$, stop at its first $B$ vertex
and begin at its last earlier $A$ vertex. The retained segment has all
internal vertices outside $A\cup B$, so every edge is weak.
Loop erasure preserves these properties. This handles arbitrary detours,
distant crossings and earlier strong-edge excursions.

Let $W$ denote existence of any finite self-avoiding all-weak open
coarse path from $A$ to $B$. Such a path starts at $(1,n)$ or $(n,1)$,
$n\geq1$, and its length $m$ is at least its distance $n+3$ to $B$.
Double-counting the corner is harmless. There are at most
$4\cdot3^{m-1}$ candidate self-avoiding paths per start. Each eligible
all-weak path has open probability $q^m$; paths containing strong edges
are again used only to enlarge the count. For $r<1$,
\begin{equation}\label{bridge}
 \PP_p(W)\leq
 2\sum_{n\geq1}\sum_{m\geq n+3}4\cdot3^{m-1}q^m
 =\frac83\,\frac{r^4}{(1-r)^2}.
\end{equation}
This sums every boundary start and every length. There is no finite-window
cutoff or assumed uniform closed barrier.

Bernoulli's inequality gives
$(20/19)^{20}\geq1+20/19>2$.
Raising to the fifth power yields $q<1/32$, so $r<3/32$.
The positive series in~\eqref{bridge} is increasing in $r$, whence
\begin{equation}\label{weakbound}
 \PP_{19/20}(W)<
 \frac83\,\frac{(3/32)^4}{(1-3/32)^2}
 =\frac{27}{107648}.
\end{equation}

The events $S_A,S_B,W^c$ depend respectively on the internal-$A$ bonds,
internal-$B$ bonds and weak-chain bonds. These are disjoint families.
Their generated sigma-fields are independent: independence for cylinder
events extends to the sigma-fields by the monotone-class argument.
All events are measurable, with survival a countable intersection of
finite-distance connection events and $W$ a countable union over finite paths.
On $S_A\cap S_B\cap W^c$, the two infinite surviving clusters cannot lie
in one actual component. Such a connection would supply a finite actual,
hence coarse, $A$-to-$B$ path and force $W$.
Consequently
\[
 \PP_{19/20}(N_\infty\geq2)
 \geq \PP(S_A)\PP(S_B)\PP(W^c)
 >\left(\frac{143}{245}\right)^2\frac{107621}{107648}
 =\frac{2200741829}{6461571200}.
\]
Subtracting $1/3$ gives $140654287/19384713600>0$.
This proves Theorem~\ref{main}. The parameter is strictly supercritical;
positive-probability nonuniqueness suffices to contradict almost-sure
uniqueness. No $0$--$1$--infinity law, deterministic cluster count,
transitivity, or uniqueness-monotonicity assertion is used.

\section{Related work, attribution and limits}
Bounded subdivision, series reduction, dual parity, path counting and
product independence are established elementary tools.
Benjamini and Schramm~\cite{bs}, Example 1 and Proposition 1,
already give a planar example with a supercritical nonuniqueness regime.
It is an enhanced binary tree roughly isometric to a hyperbolic sector,
rather than a Euclidean lattice; its stated planar model is site percolation.
Their Lemma 1 concerns uniqueness sufficiently near one.
Benjamini--Kozma~\cite[Theorem 1]{bk} give critical nonuniqueness on
a subexponential-growth product graph, without the required lattice
rough-isometry assertion.

Teixeira~\cite[Theorem 1.4]{teixeira} and
Candellero--Teixeira~\cite[Theorem 1.8]{ct} prove high-density uniqueness
under their respective local-isoperimetric/volume or rough-transitivity,
growth, isoperimetric and decoupling assumptions, with site-model distinctions.
Hutchcroft--Tointon~\cite[Theorem 1.8]{ht} give a high-density bond result
under a uniform finite-exhaustion isoperimetric hypothesis, and their
Theorem 1.10 has a separate uniform-connection premise.
A threshold below one in these theorems is not an identification with $p_c$.
The present example is compatible with a weaker near-one question.

Hil\'ario--S\'a--Sanchis--Teixeira~\cite[Section 2 and Remark 5]{hsst}
explicitly use independent coarse probabilities that are powers of $p$.
Their full-supercritical uniqueness statement concerns an \emph{annealed}
random stretched-lattice law, using its translation invariance, ergodicity
and finite energy. Those properties of the averaged law do not assert
uniqueness for every deterministic angular stretching pattern.
No comparison theorem is imported into the proof above.

The primary source question, definition, bond convention and chapter context
were read in the pinned 30 October 2013 manuscript. The published chapter
has different pagination; byte identity with it is not asserted.
Relevant statements of the comparison papers were read, not their complete
proofs. Exact versions, hashes and access limits are in the source archive.
Bounded primary-source, public-record and duplicate checks on
22 October 2026 authenticated no earlier publication of this exact result.
This is not an exhaustive novelty, exclusive-firstness, or current-openness
certificate. The earlier project \href{https://github.com/AlecKriebel/Math/pull/1348}{PR1348},
head \texttt{4222445}, is an unpublished proof source, not an earlier
published solution. Catalog 10000037/AMR-099-0037 is one result family.
The source-question DOI is not a counterexample DOI.

The source archive supplies an exact standard-library checker for finite
dual-boundary, series-connectivity, metric and rational-constant controls,
with deliberately false mutations. These are supplementary diagnostics;
the written argument supplies the unbounded paths, infinite events and
uniform metric quantifiers. We do not assert almost-sure multiplicity,
exact cluster count, critical nonuniqueness, or a simultaneous-in-$p$ law.

\paragraph{Authorship and review disclosure.}
AI tools were extensively used in all phases: planning, mathematical exploration,
proof development and reconstruction, computation, literature investigation,
writing, editing, and adversarial verification. Alec Kriebel is the named
human author. This preprint is unrefereed, without conventional independent
human peer review, journal acceptance, or formal proof-assistant certification.
Internal AI audits do not supply those certifications.
Authored prose is CC BY 4.0; authored code is MIT licensed.
Third-party bodies are cited rather than redistributed.

\begin{thebibliography}{9}\small
\bibitem{benjamini} I.~Benjamini, \emph{Coarse Geometry and Randomness},
Lecture Notes in Math.~2100, Springer, 2013.
Author manuscript, 30 October 2013, Definition 1.25 and Problem 9.6.
Published chapter: \href{https://doi.org/10.1007/978-3-319-02576-6_9}{doi:10.1007/978-3-319-02576-6\_9}.
\bibitem{bs} I.~Benjamini and O.~Schramm, \emph{Percolation beyond $\ZZ^d$,
many questions and a few answers}, Electron. Commun. Probab.~1 (1996), 71--82.
\href{https://emis.de/ft/41630}{Primary PDF}.
\bibitem{bk} I.~Benjamini and G.~Kozma, \emph{Uniqueness of percolation on
products with $\ZZ$}, ALEA~10 (2013), 15--25.
\href{https://alea.math.cnrs.fr/articles/v10/10-02.pdf}{Primary PDF}.
\bibitem{teixeira} A.~Teixeira, \emph{Percolation and local isoperimetric
inequalities}, Probab. Theory Relat. Fields~165 (2016), 963--984.
\href{https://doi.org/10.1007/s00440-015-0653-5}{doi:10.1007/s00440-015-0653-5}.
\bibitem{ct} E.~Candellero and A.~Teixeira, \emph{Percolation and isoperimetry
on roughly transitive graphs}, Ann. Inst. H.~Poincar\'e Probab. Statist.~54
(2018), 1819--1847.
\href{https://doi.org/10.1214/17-AIHP857}{doi:10.1214/17-AIHP857}.
\bibitem{ht} T.~Hutchcroft and M.~Tointon,
\emph{Non-triviality of the phase transition for percolation on finite
transitive graphs}, J.~Eur. Math. Soc.~27 (2025), 4283--4346.
\href{https://doi.org/10.4171/JEMS/1453}{doi:10.4171/JEMS/1453}.
\bibitem{hsst} M.~Hil\'ario, M.~S\'a, R.~Sanchis and A.~Teixeira,
\emph{Percolation on stretched lattices: a static renormalization approach},
Probab. Theory Relat. Fields~195 (2026), 451--496.
\href{https://doi.org/10.1007/s00440-026-01470-4}{doi:10.1007/s00440-026-01470-4}.
\end{thebibliography}
\end{document}
